% ============================================================================
%  05-tu-phan-bien — Tự phản biện: cơ chế của thất bại và bốn kỹ thuật đã đo được
%  Ngày tạo: 2026-09-09 | Sửa gần nhất: 2026-09-09 | Ôn gần nhất: chưa
%  Dependency: 02, đọc gần 04. Mức độ: CỐT LÕI.
% ============================================================================
\documentclass[../hieu-suat.tex]{subfiles}
\begin{document}
\chapter{Tự phản biện: vì sao nó thất bại, và cái gì thật sự chữa được}
\ptrtarget{05}\ptrtarget{05-tu-phan-bien}
\dependency{\ptr{02} (phản hồi hướng nhiệm vụ so với hướng bản thân), \ptr{04} (lượt viết lại đối kháng là nơi chương này được dùng).}

\begin{tomtat}
Tự tìm lỗi trong lập luận của chính mình khó không phải vì thiếu kỷ luật, mà vì hai cơ chế nhận thức có thật: bộ máy suy luận của con người vận hành tốt hơn khi \emph{đánh giá} lập luận của người khác so với khi \emph{sản xuất} lập luận cho mình, và ý tưởng đầu tiên chiếm lấy sự chú ý --- ở chuyên gia, chuyển động mắt cho thấy họ vẫn nhìn vào các đặc điểm của lời giải đầu ngay cả khi tin rằng mình đang tìm lời giải khác. Vì cơ chế nằm ở tầng chú ý, ``cố khách quan hơn'' không hiệu quả; thứ hiệu quả là các thủ tục ép sinh ra nội dung đối lập. Chương này trình bày bốn thủ tục có bằng chứng đo được và ghép chúng thành một nghi thức áp dụng trước khi tin bất kỳ kết quả nào của chính mình.
\end{tomtat}

\begin{khoigoi}
\h{Vì sao tôi bắt lỗi trong bài của người khác dễ hơn hẳn trong bài của mình?}
\h{Chuyên gia có tránh được bẫy ``ý tưởng đầu tiên'' không?}
\h{Có thủ tục nào đã đo được là làm tăng số lỗi mình tự tìm ra không, hay tất cả chỉ là lời khuyên đạo đức?}
\end{khoigoi}

\section{Cơ chế thứ nhất: suy luận được thiết kế để tranh luận}

Mercier và Sperber \cite{mercier2011} lập luận rằng chức năng tiến hoá của khả năng suy luận là \emph{tạo và đánh giá lập luận trong giao tiếp}, chứ không phải để một cá nhân tự tìm ra chân lý. Từ đó suy ra một dự đoán kiểm được: con người thiên vị khi \emph{sản xuất} lập luận cho phía mình, nhưng khách quan hơn nhiều khi \emph{đánh giá} lập luận của người khác \nT{}. Bằng chứng họ tập hợp gồm hiện tượng quen thuộc trong các bài toán suy luận chuẩn: tỉ lệ giải đúng của cá nhân rất thấp --- cỡ 10\% --- nhưng khi cùng bài toán đó được thảo luận trong nhóm, tỉ lệ lên khoảng 80\% \nD{}.

Hệ quả trực tiếp và hơi khó chịu \nM{}: ``tự phản biện'' theo nghĩa đen là yêu cầu bộ máy sản xuất tự đóng vai bộ máy đánh giá. Nó không hỏng hoàn toàn, nhưng nó là chế độ vận hành kém tự nhiên nhất, nên phải được chống đỡ bằng thủ tục bên ngoài chứ không bằng quyết tâm.

\section{Cơ chế thứ hai: ý tưởng đầu tiên chiếm lấy sự chú ý}

Bilalić, McLeod và Gobet \cite{bilalic2008} đặt kỳ thủ trước một thế cờ có hai lời giải: một lời giải quen thuộc, dài; và một lời giải khác ngắn hơn, tốt hơn. Sau khi tìm ra lời giải quen, các kỳ thủ nói rằng họ đang tìm lời giải tốt hơn --- nhưng chuyển động mắt cho thấy họ vẫn tiếp tục nhìn vào những ô liên quan tới lời giải đã nghĩ ra \nD{}. Cơ chế mà các tác giả rút ra: ý tưởng đầu tiên định hướng chú ý về phía thông tin phù hợp với nó và kéo chú ý ra khỏi thông tin không phù hợp, một cách vô thức \nT{}.

\begin{cambay}
Điều này phá hỏng phép tự kiểm phổ biến nhất: ``tôi đã cân nhắc các khả năng khác rồi''. Các kỳ thủ trong thí nghiệm cũng tin đúng như vậy, và họ sai --- không phải nói dối, mà là không quan sát được chính mình. Cảm giác đã cân nhắc không phải bằng chứng của việc đã cân nhắc \nM{}.
\end{cambay}

\section{Bốn thủ tục có bằng chứng}

\textbf{Một --- xét mặt đối lập.} Lord, Lepper và Preston \cite{lord1984} so sánh hai chỉ dẫn: ``hãy cố công bằng và không thiên vị'' so với ``hãy tự hỏi bạn sẽ đánh giá thế nào nếu bằng chứng chỉ theo hướng ngược lại''. Chỉ dẫn thứ hai làm giảm thiên lệch mạnh hơn hẳn chỉ dẫn thứ nhất \nD{}. Bài học: lời kêu gọi khách quan không có tác dụng; một thao tác cụ thể sinh ra nội dung đối lập thì có.

\textbf{Hai --- hồi tưởng từ tương lai (premortem).} Mitchell, Russo và Pennington \cite{mitchell1989} cho thấy khi người ta được yêu cầu tưởng tượng một sự kiện tương lai \emph{đã xảy ra rồi} và giải thích vì sao, họ sinh ra nhiều hơn khoảng 30\% số lý do so với khi được hỏi sự kiện đó \emph{có thể} xảy ra vì sao \nD{}; điều tạo ra khác biệt là tính chắc chắn của cách đặt câu hỏi. Klein \cite{klein2007} biến điều đó thành thủ tục nhóm: trước khi bắt đầu, giả định dự án đã thất bại thảm hại, và mỗi người viết ra vì sao \nT{}.

\textbf{Ba --- huấn luyện hiệu chỉnh và ghi lại dự đoán bằng số.} Trong giải đấu dự báo của Mellers và cộng sự \cite{mellers2015}, nhóm dự báo giỏi nhất được xác định và bồi dưỡng có hệ thống; mô hình tốt nhất kết hợp đặc điểm cá nhân, tình huống và hành vi cho tương quan bội 0,64 với độ chính xác \nD{}. Điều đáng mang sang nghiên cứu \nM{} không phải là con số đó mà là hình thức: dự đoán được ghi bằng \emph{xác suất}, và điểm số được tính về sau, nên người dự báo có một vòng phản hồi định lượng --- đúng thứ mà chương \ptr{02} nói môi trường nghiên cứu không có.

\textbf{Bốn --- mượn bộ máy đánh giá của người khác.} Dunbar \cite{dunbar1997} quan sát trực tiếp bốn phòng thí nghiệm và đo được điều mà lý thuyết của Mercier và Sperber dự đoán: trong một buổi họp, khi một nghiên cứu viên trình bày 5 nhóm kết quả kèm 11 suy luận, các thành viên khác đã giới hạn, mở rộng, thay thế hoặc loại bỏ \textbf{7 trong 11} suy luận đó \nD{}. Ông cũng đếm mức chú ý của nhóm dành cho các phát hiện bất ngờ: 176 lượt tương tác cho các phát hiện bất ngờ so với 23 lượt cho các phát hiện đúng như dự đoán \nD{}.

\begin{trucgiac}
Ba thủ tục đầu là những cách \emph{giả lập} người phản biện bên ngoài khi không có ai. Thủ tục thứ tư là có người thật. Thứ tự ưu tiên vì thế rõ ràng \nM{}: nếu có thể lấy được một người đọc thật --- kể cả người không cùng chuyên môn hẹp --- thì đó là lựa chọn tốt nhất, và ba thủ tục kia dùng cho những ngày không có ai.
\end{trucgiac}

\section{Chuẩn mực gốc: nghiêng hẳn về phía bất lợi cho mình}

Feynman \cite{feynman1974} đặt tên cho chuẩn mực này trong bài nói năm 1974: nguyên tắc thứ nhất là không được tự lừa mình, mà mình lại là người dễ tự lừa nhất; và điều đó đòi hỏi một kiểu chính trực ``nghiêng hẳn về phía ngược lại'' --- báo cáo mọi thứ có thể làm cho kết quả của mình sai, chứ không chỉ những gì ủng hộ nó \nT{}. Đây là bài luận, không phải nghiên cứu thực nghiệm; giá trị của nó là ở chỗ nó phát biểu tiêu chuẩn, còn cơ chế vì sao tiêu chuẩn đó khó đạt thì hai mục đầu chương đã nói.

\section{Nghi thức trước khi tin một kết quả của chính mình}

Ghép lại thành một thủ tục dùng được, tốn khoảng nửa giờ, chạy \emph{trước} khi báo cáo bất kỳ kết quả nào \nM{}:

\begin{enumerate}
\item \textbf{Viết ra dự đoán trước khi nhìn số} --- kèm xác suất. Không có bước này thì mọi kết quả đều sẽ trông hợp lý sau khi đã biết.
\item \textbf{Premortem:} ``sáu tháng nữa kết quả này hoá ra sai. Vì sao?'' Viết ít nhất năm lý do, chắc chắn ở dạng đã xảy ra chứ không phải có thể xảy ra \cite{mitchell1989,klein2007}.
\item \textbf{Xét mặt đối lập:} nếu con số ngược lại, tôi sẽ giải thích thế nào? Nếu tôi giải thích được cả hai chiều dễ như nhau thì thí nghiệm này không phân biệt được giả thuyết nào cả \cite{lord1984} \nM{}.
\item \textbf{Kiểm tra giả thuyết cạnh tranh tầm thường:} nhiễu giữa các lần chạy, khác biệt về công sức tinh chỉnh giữa hai nhánh, rò rỉ dữ liệu, một lỗi cài đặt có lợi cho phía mình. Chương \ptr{08} liệt kê danh sách đầy đủ cho CV/ML.
\item \textbf{Đưa cho một người thật} với yêu cầu cụ thể: ``tìm giúp tôi lý do kết quả này sai'', chứ không phải ``xem giúp có ổn không'' \cite{mercier2011,dunbar1997}.
\item \textbf{Ghi vào sổ dự đoán} kết quả cuối so với dự đoán ban đầu. Sau vài chục lần, sổ này cho biết trực giác của tôi đang hiệu chỉnh tốt tới đâu \cite{mellers2015}.
\end{enumerate}

\begin{cannho}
``Tôi đã kiểm tra kỹ rồi'' không phải một bước kiểm tra. Bốn cơ chế bảo vệ duy nhất đã được đo là: dự đoán viết trước, sinh nội dung đối lập một cách bắt buộc, giả định thất bại đã xảy ra, và người thứ hai đọc với nhiệm vụ bác bỏ.
\end{cannho}

\begin{traloi}
\tl{Vì bộ máy suy luận vận hành ở chế độ đánh giá tốt hơn ở chế độ sản xuất \cite{mercier2011}; tự phản biện là bắt cùng một bộ máy làm cả hai vai.}
\tl{Không. Kỳ thủ mạnh vẫn bị lời giải đầu tiên giữ lấy sự chú ý, và họ không nhận ra điều đó --- chuyển động mắt mâu thuẫn với lời họ nói \cite{bilalic2008}.}
\tl{Có: xét mặt đối lập mạnh hơn lời kêu gọi khách quan \cite{lord1984}; giả định sự kiện đã xảy ra sinh thêm khoảng 30\% lý do \cite{mitchell1989}; ghi dự đoán bằng xác suất cho vòng phản hồi định lượng \cite{mellers2015}; và nhóm thật sự chặn được suy luận sai của cá nhân \cite{dunbar1997}.}
\end{traloi}

\begin{tukiem}
\item Nêu hai cơ chế khiến tự phản biện thất bại, và giải thích vì sao cả hai đều không chữa được bằng cách ``cố gắng khách quan hơn''.
\item Vì sao dạng câu hỏi ``vì sao nó \emph{đã} thất bại'' sinh ra nhiều lý do hơn ``vì sao nó \emph{có thể} thất bại''?
\item Lấy kết quả gần nhất của bạn và chạy bước 3 của nghi thức: nếu con số ngược lại, bạn sẽ giải thích thế nào? Bạn giải thích có dễ không --- và điều đó nói gì về thí nghiệm?
\item Trong dữ liệu của Dunbar, nhóm đã sửa 7 trong 11 suy luận của người trình bày. Bạn có một cấu trúc tương đương trong tuần làm việc của mình không? Nếu không, cái gì thay thế được?
\end{tukiem}

\begin{onbai}
\ar[3]{Sản xuất vs đánh giá lập luận}{Con người khách quan hơn khi đánh giá lập luận người khác; nhóm 80\% so với cá nhân 10\% ở bài suy luận chuẩn \cite{mercier2011}.}
\ar[3]{Einstellung}{Ý tưởng đầu tiên chiếm chú ý; chuyên gia tin mình đang tìm cái khác nhưng mắt vẫn nhìn cái cũ \cite{bilalic2008}.}
\ar[3]{Xét mặt đối lập}{Sinh nội dung đối lập cụ thể mạnh hơn hẳn chỉ dẫn ``hãy công bằng'' \cite{lord1984}.}
\ar[3]{Premortem}{Giả định thất bại \emph{đã} xảy ra \(\Rightarrow\) khoảng +30\% lý do \cite{mitchell1989,klein2007}.}
\ar[2]{Sổ dự đoán xác suất}{Cách duy nhất biến trực giác thành thứ hiệu chỉnh được \cite{mellers2015}.}
\ar[2]{Nhóm sửa suy luận}{7/11 suy luận bị nhóm sửa; 176 vs 23 lượt tương tác cho phát hiện bất ngờ \cite{dunbar1997}.}
\ar[1]{Chính trực kiểu Feynman}{Báo cáo mọi thứ có thể làm mình sai, không chỉ thứ ủng hộ mình \cite{feynman1974}.}
\end{onbai}

\end{document}
